\section{Related Work}

Existing research on generative AI security mainly focuses on the safety of model outputs. This line of work studies various attack methods \citep{li2024badedit, li2025one} and corresponding defenses \citep{zhang2025exploiting}. However, these attacks lack a unified definition of security. A representative example is the jailbreak attack \citep{mehrotra2024tree, zhang2025jbshield, yu2024don, wang-etal-2025-diffusionattacker}. This attack aims to induce models to generate unsafe content. Such content covers multiple categories \citep{guo2025deepseeknature}, including discrimination, illegal behavior, harmful behavior, and ethical issues. Output level unsafety is inherently subjective. This subjectivity prevents a deterministic definition under cryptographic theory. Current evaluation of output safety largely relies on LLM-as-a-judge \citep{li-etal-2025-generation}. This approach depends on specialized evaluator models trained for safety analysis, such as Qwen3Guard \citep{zhao2025qwen3guard}. These models perform semantic analysis on generated content. However, evaluator models are not interpretable. Their correctness in safety assessment cannot be formally proven. Defenses against these attacks also lack formal guarantees. As a result, research on generative model security lacks a mature theoretical foundation. This absence of theory leads to inconsistent standards across existing attack and defense studies. 
% In contrast, our work provides a theoretical guarantee, ensuring that the delivered response admits an upper bound on risk.
